% GENERATED FILE. Edit canonical lessons, then run npm run generate:books.
% canonical-source-hash: fnv1a64:33ef66f21836e7c6
% canonical-lessons: FA-C380-rezume, FA-C380-qarardad, FA-C380-tarfi, FA-C380-karamuz, FA-C380-karamuzi

\chapter{Applying for work}
\label{ch:fa-applying-for-work}
\hlchaptermodality{380}

\begin{chapteropening}
\textbf{By the end of this chapter:} \emph{I can talk about sending a CV, signing a contract, getting a promotion and doing an internship.}

Complete the last lesson of chapter 380: I can talk about sending a CV, signing a contract, getting a promotion and doing an internship.
\end{chapteropening}

\section[\texorpdfstring{\fa{رزومه} (rezume)}{رزومه (rezume)}]{\fa{رزومه} (rezume) — a CV, a résumé}

\label{lesson:FA-C380-rezume}



\begin{quote}
Before the new one: say the Persian for unemployment, then the Persian for retirement.
\end{quote}

\subsection*{You'll want to know: \fa{رزومه}}
\textbf{\fa{رزومه}} (\emph{rezume}) — \textquotedblleft{}a CV, a résumé\textquotedblright{}.

\begin{quote}
\fa{رزومه} \fa{خودم} \fa{را} \fa{برای} \fa{سه} \fa{شرکت} \fa{فرستادم}.
\end{quote}

\begin{quote}
\emph{rezume-ye khodam râ barâye se sherkat ferestâdam.}
\end{quote}

\begin{quote}
I sent my CV to three companies.
\end{quote}

A loan from French. A more Persian phrase is \emph{sâbeqe-ye kâri}, a work record.

\begin{cousinweb}
\textbf{\fa{رزومه}} was borrowed from French \emph{résumé} \textquotedblleft{}a summary\textquotedblright{}, the participle of \emph{résumer}, from Latin \emph{resumere} \textquotedblleft{}to take up again\textquotedblright{} (Lewis \& Short): a CV sums up your working life.
\end{cousinweb}

\begin{grammarlens}[title={What you've built}]
One more word for telling, explaining and describing at B1.
\end{grammarlens}

\subsection*{Your turn}
\begin{itemize}
\raggedright
  \item \emph{Say it:} \emph{rezume}
  \item \emph{Say it:} \emph{rezume}, once more
  \item \emph{Say it:} a sentence of your own with \emph{rezume}
  \item \emph{Recall:} say the Persian for unemployment, then the Persian for retirement, then say \emph{rezume} again
\end{itemize}

\begin{tcolorbox}[breakable,colback=teal!4,colframe=teal!35!black,title={Before you move on}]
What does \fa{رزومه} mean? (\textquotedblleft{}A CV, a résumé\textquotedblright{}.) Say it once more.
\end{tcolorbox}
\section[\texorpdfstring{\fa{قرارداد} (qarârdâd)}{قرارداد (qarârdâd)}]{\fa{قرارداد} (qarârdâd) — a contract}

\label{lesson:FA-C380-qarardad}



\begin{quote}
Before the new one: say the Persian for retirement, then the Persian for a CV.
\end{quote}

\subsection*{You'll want to know: \fa{قرارداد}}
\textbf{\fa{قرارداد}} (\emph{qarârdâd}) — \textquotedblleft{}a contract\textquotedblright{}.

\begin{quote}
\fa{او} \fa{یک} \fa{قرارداد} \fa{یک} \fa{ساله} \fa{امضا} \fa{کرد}.
\end{quote}

\begin{quote}
\emph{u yek qarârdâd-e yek sâle emzâ kard.}
\end{quote}

\begin{quote}
He signed a one-year contract.
\end{quote}

\emph{Qarâr} \textquotedblleft{}agreement\textquotedblright{} + \emph{dâd} \textquotedblleft{}gave\textquotedblright{}. \emph{Qarârdâd bastan} — to sign a contract.

\begin{cousinweb}
\textbf{\fa{قرارداد}} joins \emph{qarâr}, from Arabic \emph{qarār} \textquotedblleft{}a settled decision\textquotedblright{}, to \emph{dâd}, the past stem of \emph{dâdan} \textquotedblleft{}to give\textquotedblright{}, which EtymDB traces to PIE \emph{deh\textsubscript{3}-}: an agreement that has been given.
\end{cousinweb}

\begin{grammarlens}[title={What you've built}]
One more word for telling, explaining and describing at B1.
\end{grammarlens}

\subsection*{Your turn}
\begin{itemize}
\raggedright
  \item \emph{Say it:} \emph{qarârdâd}
  \item \emph{Say it:} \emph{qarârdâd}, once more
  \item \emph{Say it:} a sentence of your own with \emph{qarârdâd}
  \item \emph{Recall:} say the Persian for retirement, then the Persian for a CV, then say \emph{qarârdâd} again
\end{itemize}

\begin{tcolorbox}[breakable,colback=teal!4,colframe=teal!35!black,title={Before you move on}]
What does \fa{قرارداد} mean? (\textquotedblleft{}A contract\textquotedblright{}.) Say it once more.
\end{tcolorbox}
\section[\texorpdfstring{\fa{ترفیع} (tarfi')}{ترفیع (tarfi')}]{\fa{ترفیع} (tarfi') — a promotion}

\label{lesson:FA-C380-tarfi}



\begin{quote}
Before the new one: say the Persian for a CV, then the Persian for a contract.
\end{quote}

\subsection*{You'll want to know: \fa{ترفیع}}
\textbf{\fa{ترفیع}} (\emph{tarfi'}) — \textquotedblleft{}a promotion\textquotedblright{}.

\begin{quote}
\fa{خواهرم} \fa{امسال} \fa{ترفیع} \fa{گرفت}.
\end{quote}

\begin{quote}
\emph{khâharam emsâl tarfi' gereft.}
\end{quote}

\begin{quote}
My sister got a promotion this year.
\end{quote}

\emph{Tarfi' gereftan} — to be promoted.

\begin{cousinweb}
\textbf{\fa{ترفیع}} is usually traced to Arabic \emph{tarfī'} \textquotedblleft{}raising, lifting up\textquotedblright{}: a promotion lifts you to a higher post.
\end{cousinweb}

\begin{grammarlens}[title={What you've built}]
One more word for telling, explaining and describing at B1.
\end{grammarlens}

\subsection*{Your turn}
\begin{itemize}
\raggedright
  \item \emph{Say it:} \emph{tarfi'}
  \item \emph{Say it:} \emph{tarfi'}, once more
  \item \emph{Say it:} a sentence of your own with \emph{tarfi'}
  \item \emph{Recall:} say the Persian for a CV, then the Persian for a contract, then say \emph{tarfi'} again
\end{itemize}

\begin{tcolorbox}[breakable,colback=teal!4,colframe=teal!35!black,title={Before you move on}]
What does \fa{ترفیع} mean? (\textquotedblleft{}A promotion\textquotedblright{}.) Say it once more.
\end{tcolorbox}
\section[\texorpdfstring{\fa{کارآموز} (kârâmuz)}{کارآموز (kârâmuz)}]{\fa{کارآموز} (kârâmuz) — a trainee, an intern}

\label{lesson:FA-C380-karamuz}



\begin{quote}
Before the new one: say the Persian for a contract, then the Persian for a promotion.
\end{quote}

\subsection*{You'll want to know: \fa{کارآموز}}
\textbf{\fa{کارآموز}} (\emph{kârâmuz}) — \textquotedblleft{}a trainee, an intern\textquotedblright{}.

\begin{quote}
\fa{سه} \fa{کارآموز} \fa{در} \fa{دفتر} \fa{ما} \fa{کار} \fa{میکنند}.
\end{quote}

\begin{quote}
\emph{se kârâmuz dar daftar-e mâ kâr mikonand.}
\end{quote}

\begin{quote}
Three trainees work in our office.
\end{quote}

\emph{Kâr} \textquotedblleft{}work\textquotedblright{} + \emph{âmuz}, the present stem of \emph{âmukhtan} \textquotedblleft{}to learn\textquotedblright{}. Plural \emph{kârâmuzân}.

\begin{cousinweb}
\textbf{\fa{کارآموز}} is a native Persian compound built from \emph{kâr} \textquotedblleft{}work\textquotedblright{} + \emph{âmuz-}, the present stem of \emph{âmukhtan} \textquotedblleft{}to learn, to teach\textquotedblright{} (EtymDB traces it to Middle Persian). A trainee is a work-learner.
\end{cousinweb}

\begin{grammarlens}[title={What you've built}]
One more word for telling, explaining and describing at B1.
\end{grammarlens}

\subsection*{Your turn}
\begin{itemize}
\raggedright
  \item \emph{Say it:} \emph{kârâmuz}
  \item \emph{Say it:} \emph{kârâmuz}, once more
  \item \emph{Say it:} a sentence of your own with \emph{kârâmuz}
  \item \emph{Recall:} say the Persian for a contract, then the Persian for a promotion, then say \emph{kârâmuz} again
\end{itemize}

\begin{tcolorbox}[breakable,colback=teal!4,colframe=teal!35!black,title={Before you move on}]
What does \fa{کارآموز} mean? (\textquotedblleft{}A trainee, an intern\textquotedblright{}.) Say it once more.
\end{tcolorbox}
\section[\texorpdfstring{\fa{کارآموزی} (kârâmuzi)}{کارآموزی (kârâmuzi)}]{\fa{کارآموزی} (kârâmuzi) — an internship, training}

\label{lesson:FA-C380-karamuzi}



\begin{quote}
Before the new one: say the Persian for a promotion, then the Persian for a trainee.
\end{quote}

\subsection*{You'll want to know: \fa{کارآموزی}}
\textbf{\fa{کارآموزی}} (\emph{kârâmuzi}) — \textquotedblleft{}an internship, training\textquotedblright{}.

\begin{quote}
\fa{او} \fa{تابستان} \fa{دوره} \fa{کارآموزی} \fa{را} \fa{در} \fa{یک} \fa{بیمارستان} \fa{گذراند}.
\end{quote}

\begin{quote}
\emph{u tâbestân dowre-ye kârâmuzi râ dar yek bimârestân gozarând.}
\end{quote}

\begin{quote}
In the summer he did his internship in a hospital.
\end{quote}

\emph{Kârâmuz} + the noun suffix \emph{-i}. \emph{Dowre-ye kârâmuzi} — a training period.

\begin{cousinweb}
\textbf{\fa{کارآموزی}} is native Persian, built from \emph{kârâmuz} \textquotedblleft{}trainee\textquotedblright{}, from \emph{kâr} \textquotedblleft{}work\textquotedblright{} and \emph{âmuz-} \textquotedblleft{}learn\textquotedblright{}, with the noun suffix \emph{-i}.
\end{cousinweb}

\begin{grammarlens}[title={What you've built}]
One more word for telling, explaining and describing at B1.
\end{grammarlens}

\subsection*{Your turn}
\begin{itemize}
\raggedright
  \item \emph{Say it:} \emph{kârâmuzi}
  \item \emph{Say it:} \emph{kârâmuzi}, once more
  \item \emph{Say it:} a sentence of your own with \emph{kârâmuzi}
  \item \emph{Recall:} say the Persian for a promotion, then the Persian for a trainee, then say \emph{kârâmuzi} again
\end{itemize}

\begin{tcolorbox}[breakable,colback=teal!4,colframe=teal!35!black,title={Before you move on}]
What does \fa{کارآموزی} mean? (\textquotedblleft{}An internship, training\textquotedblright{}.) Say it once more.
\end{tcolorbox}
